\documentclass{article}
\usepackage[utf8]{inputenc}
\usepackage[T1]{fontenc}
\usepackage{geometry}
\geometry{
    a4paper,
    left=2.5cm,
    right=2.5cm,
    top=2.5cm,
    bottom=2.5cm
}
\usepackage{graphicx}
\usepackage{float}
\usepackage{hyperref}
\usepackage[english]{babel}
\usepackage{amsmath}
\usepackage{amssymb}
\usepackage{array}
\usepackage{booktabs}
\usepackage{caption}
\usepackage{subcaption}
\usepackage{enumitem}
\usepackage{fancyhdr}
\usepackage{titlesec}
\usepackage{url}
\usepackage{textcomp}
\usepackage{xcolor}
\usepackage{listings}
\usepackage{multirow}
\usepackage{longtable}
\usepackage{tabularx}
\usepackage{times}

% Set up page style
\pagestyle{fancy}
\fancyhf{}
\rfoot{\thepage}
\renewcommand{\headrulewidth}{0pt}

% Title formatting
\titleformat{\chapter}[display]{\normalfont\bfseries}{\chaptertitle\thechapter}{0pt}{\Huge}
\titlespacing*{\chapter}{0pt}{50pt}{40pt}

% Document info
\title{Template Document}
\author{Academic Support Agent}
\date{\today}

\begin{document}
\maketitle

%----------------------------------------------------------------------------------------
%	TITLE PAGE
%----------------------------------------------------------------------------------------

\begin{frame}
\titlepage
\end{frame}				% 

\begin{frame}{ } % Table of contents slide, comment this block out to remove it
%


\end{frame}

\AtBeginSection[]
{
  \begin{frame}
    \frametitle{ }
    
  \end{frame}
}


%----------------------------------------------------------------------------------------
%	PRESENTATION SLIDES
%----------------------------------------------------------------------------------------

%------------------------------------------------
\section{ }% Sections can be created in order to organize your presentation into discrete blocks, all sections and subsections are automatically printed in the table of contents as an overview of the talk
%------------------------------------------------

\begin{frame}
  \frametitle{ }

  \begin{itemize}
    \item { }
	    \begin{itemize}
	    	\item    Overleaf
	    	\item   \XeLaTeX  
	    \end{itemize}
    \item   \LaTeX   Beamer   
    
    \item   $\sin^2 \theta + \cos^2 \theta = 1$
    \item {  \begin{equation}
	    y_{1}=\int \sin x\, {\rm d}x
    \end{equation}	 }   
  \end{itemize}
\end{frame}
%------------------------------------------------

%------------------------------------------------
\section{ }
%------------------------------------------------

\begin{frame}
  \frametitle{ }
	\begin{block}{Slides with \LaTeX}
	    Beamer offers a lot of functions to create nice slides using \LaTeX.
	  \end{block}
	
	  \begin{block}{The basis}
	     
	    \begin{itemize}
	      \item split
	      \item whale
	      \item rounded
	      \item orchid
	    \end{itemize}
	  \end{block}
\end{frame}
%------------------------------------------------
\begin{frame}
  \frametitle{ }
	 \begin{enumerate}
	    \item This just shows the effect of the style
	    \item It is not a Beamer tutorial
	    \item Read the Beamer manual for more help
	    \item Contact me only concerning the style file
	  \end{enumerate}
\end{frame}
%------------------------------------------------
\begin{frame}{ }
\begin{columns}
\begin{column}{0.5\textwidth}
\begin{block}{block1}
%\footnotesize{
\begin{enumerate}
        \item item1
        \begin{itemize}
            \item  item1.1
        \end{itemize}
        \item item2
        \begin{itemize}
            \item item2.1
        \end{itemize}
        \item item3
        \begin{itemize}
            \item item3.1
        \end{itemize}
\end{enumerate}    
%}    
\end{block}
\end{column}
\begin{column}{0.5\textwidth}
\begin{block}{block2}
%\footnotesize{
\begin{enumerate}
        \item item1
        \begin{itemize}
            \item  item1.1
        \end{itemize}
        \item item2
        \begin{itemize}
            \item item2.1
        \end{itemize}
        \item item3
        \begin{itemize}
            \item item3.1
        \end{itemize}
\end{enumerate}    
%}    
\end{block}
\end{column}
\end{columns}
\end{frame}
%------------------------------------------------

\begin{frame}{ }
    \begin{minipage}{0.4\linewidth}
\begin{figure}[h]
            \centering
            \includegraphics[height=.4\textheight]{res/GDUT-gate.png}
            \caption{GDUT}
 \end{figure}
    \end{minipage}\hspace{0.3cm}
    \begin{minipage}{0.45\linewidth}
    \begin{figure}[h]
            \centering
            \includegraphics[height=.4\textheight]{res/GDUT-library.png}
            \caption{GDUT}
 \end{figure}   
    \end{minipage}%\hspace{0.1cm}
\end{frame}
%------------------------------------------------

\begin{frame}[t]{ }
\begin{small}
\begin{block}{ }
\begin{minipage}{0.5\linewidth}
\begin{figure}[h]
            \centering
            \includegraphics[height=.2\textheight]{res/GDUT-logo.png}
            \caption{GDUT}
 \end{figure}
    \end{minipage}\hspace{0.1cm}
    \begin{minipage}{0.45\linewidth}
    \begin{figure}[h]
            \centering
            \includegraphics[height=.2\textheight]{res/GDUT-logo_black.png}
            \caption{GDUT}
 \end{figure}   
    \end{minipage}\hspace{0.1cm}
    %\medskip
        \begin{minipage}{0.5\linewidth}
    \begin{figure}[h]
            \centering
            \includegraphics[height=.2\textheight]{res/GDUT-logo_white.png}
            \caption{GDUT}
 \end{figure}   
    \end{minipage}\hspace{0.4cm}
\begin{minipage}{0.4\linewidth}
\begin{block}{}
\begin{itemize}
\scriptsize{
\item ****************
\item ****************
\item ****************
}
\end{itemize}
\end{block}
\end{minipage}
\end{block}
\end{small}
\end{frame}
%------------------------------------------------


%------------------------------------------------
\section{ }
%------------------------------------------------
\begin{frame}
\subsection{  \frametitle{ }}

  \begin{itemize}
    \item Easy to use
    \item Good results
  \end{itemize}
\end{frame}
%------------------------------------------------

%------------------------------------------------
\section{ }
%------------------------------------------------
\begin{frame}
\frametitle{ }
\begin{block}{}
\begin{minipage}{1\linewidth}
\scriptsize{
\begin{thebibliography}{99} % Beamer does not support BibTeX so references must be inserted manually as below
\bibitem{1} R. Sun, Y. Wang, L. Lyu, N. Cheng, S. Zhang, T. Yang, and X.Shen, Delay-oriented caching strategies in d2d mobile networks,  \emph{IEEE Trans. Veh. Technol.}, vol. 69, no. 8, pp. 8529--8541, Aug. 2020.
\item Z. Su, Y. Hui, Q. Xu, T. Yang, J. Liu, and Y. Jia, An edge caching scheme to distribute content in vehicular networks,  \emph{IEEE Trans. Veh. Technol.}, vol. 67, no. 6, pp. 5346--5356, Jun. 2018.
\item Q. Xu, Z. Su, Y. Wang, and K. Zhang, Secure edge caching for layered multimedia contents in heterogeneous networks,  in \emph{Proc. IEEE Global Commun. Conf.}, Waikoloa, HI, USA, Dec. 2019, pp. 1--6.
\item B. Hu, L. Fang, X. Cheng, and L. Yang, In-vehicle caching (iv-cache) via dynamic distributed storage relay ($d^2$sr) in vehicular networks,  \emph{IEEE Trans. Veh. Technol.}, vol. 68, no. 1, pp. 843--855, Jan. 2019.
\item C. Liu, K. Liu, S. Guo, R. Xie, V. C. S. Lee, and S. H. Son, Adaptive offloading for time-critical tasks in heterogeneous internet of vehicles,  \emph{IEEE Internet Things J.}, vol. 7, no. 9, pp. 7999--8011, Sep. 2020.
\item J. Chen, H. Wu, P. Yang, F. Lyu, and X. Shen, Cooperative edge caching with location-based and popular contents for vehicular networks,  \emph{IEEE Trans. Veh. Technol.}, vol. 69, no. 9, pp. 10291-10305, Jun. 2020.
\end{thebibliography}
}
\end{minipage}
\end{block}
\end{frame}
%------------------------------------------------


%------------------------------------------------
\section{}

\begin{frame}
\begin{block}{Ending}
\Huge{\centerline{\emph{Thanks for Your Attention!}}}
\Huge{\centerline{\emph{Q \& A ?}}}
\end{block}

\end{frame}
%------------------------------------------------
\end{document}